% Compile with: latexmk -pdf main
\documentclass[11pt]{article}

\usepackage[margin=1in]{geometry}
\usepackage{amsmath,amssymb,amsthm}
\usepackage[T1]{fontenc}
\usepackage[utf8]{inputenc}
\usepackage{lmodern}
\usepackage{newunicodechar}
\usepackage{fancyvrb}
\usepackage{xcolor}
\usepackage[hidelinks]{hyperref}
\usepackage{booktabs}

% Unicode characters appearing in Lean code
\newunicodechar{‖}{\ensuremath{\|}}
\newunicodechar{•}{\ensuremath{\bullet}}
\newunicodechar{₀}{\ensuremath{{}_0}}
\newunicodechar{₃}{\ensuremath{{}_3}}
\newunicodechar{ℝ}{\ensuremath{\mathbb{R}}}
\newunicodechar{→}{\ensuremath{\to}}
\newunicodechar{↔}{\ensuremath{\leftrightarrow}}
\newunicodechar{∀}{\ensuremath{\forall}}
\newunicodechar{≠}{\ensuremath{\neq}}
\newunicodechar{⟪}{\ensuremath{\langle\!\langle}}
\newunicodechar{⟫}{\ensuremath{\rangle\!\rangle}}
\newunicodechar{⨯}{\ensuremath{\times}}

\DefineVerbatimEnvironment{lean}{Verbatim}{fontsize=\small,frame=leftline,framerule=0.6pt,
  rulecolor=\color{gray!60},xleftmargin=1em}

\newtheorem{theorem}{Theorem}
\newtheorem{lemma}[theorem]{Lemma}
\theoremstyle{definition}
\newtheorem{definition}[theorem]{Definition}
\theoremstyle{remark}
\newtheorem{remark}[theorem]{Remark}

\newcommand{\vA}{\mathbf{A}}
\newcommand{\vL}{\mathbf{L}}
\newcommand{\vp}{\mathbf{p}}
\newcommand{\vq}{\mathbf{q}}
\newcommand{\ip}[2]{\langle #1, #2\rangle}

\title{A Formal Proof of the Conservation of the\\ Laplace--Runge--Lenz Vector and Kepler's First Law\\ in Lean~4}
\author{Nichlas Lofdahl \and Claude (Anthropic)\thanks{AI system (model Claude Opus 5.5), developed by Anthropic.
Claude wrote the Lean proofs and drafted this manuscript in collaboration with the first author.}}
\date{October 2026}

\begin{document}
\sloppy
\maketitle

\begin{abstract}
We present a machine-checked formalization, in the Lean~4 proof assistant on top of the Physlib
library, of the classical Kepler problem: a particle in the inverse-square potential
$V(r)=-k/r$. We define the Laplace--Runge--Lenz (LRL) vector on phase space and prove that its
Poisson bracket with the Hamiltonian vanishes away from the centre of force. We prove the
algebraic identities $\ip{\vA}{\vL}=0$, $\ip{\vA}{\vq}=\|\vL\|^2-mk\|\vq\|$ and
$\|\vA\|^2=(mk)^2+2mH\|\vL\|^2$, and derive that bound orbits ($H<0$) have eccentricity less than
one. Finally, we lift these phase-space statements to trajectories: along every differentiable
solution of Hamilton's equations avoiding the origin, the LRL vector and $\|\vL\|^2$ are constant,
so the position satisfies the polar equation of a conic with a focus at the centre of force
(Kepler's first law). The development is about 300 lines, uses no axioms beyond Lean's standard
three, and passes all of Physlib's contribution linters.
\end{abstract}

\section{Introduction}

Formal verification of physics has gained momentum with the growth of Lean's
mathematical library Mathlib~\cite{mathlib2020} and the physics library
Physlib (formerly HepLean/PhysLean)~\cite{toobysmith2024heplean}. Physlib now contains, among much else,
a quantum-mechanical treatment of the hydrogen atom including commutation relations of the
quantum LRL operator, and, as of October~2026, basic infrastructure for classical central-force
problems: the Hamiltonian $H(\vp,\vq)=\ip{\vp}{\vp}/2m+V(\|\vq\|)$, its gradients, and the
conservation of angular momentum in Poisson-bracket form.

What was missing is the defining classical result of the Kepler problem: the existence of the
additional conserved vector $\vA$ that explains why Kepler orbits close, and the derivation of
the conic orbit equation from it. Before starting we searched Physlib's sources, its open
pull requests and issues, and the web; we are not aware of an existing formalization of the
classical LRL vector or of the conic orbit theorem in Lean. This note fills that gap.

\paragraph{Contributions.}
\begin{itemize}
  \item A definition of the Kepler problem as an instance of Physlib's \texttt{CentralForce}
  structure, and of the LRL vector as a function on phase space (\S\ref{sec:defs}).
  \item Algebraic identities: the triple-product expansion of $\vA$, orthogonality to $\vL$, the
  orbit equation, the norm identity, and the eccentricity--energy criterion (\S\ref{sec:algebra}).
  \item Conservation: explicit phase-space gradients of each component of $\vA$ and the
  vanishing of $\{A_i,H\}$ (\S\ref{sec:conservation}).
  \item Trajectory results: a general lemma that a Poisson-commuting observable is constant along
  solutions of Hamilton's equations, and Kepler's first law (\S\ref{sec:trajectory}).
\end{itemize}

The source is available at \url{https://github.com/Crotalus/kepler-lrl-lean}; it builds against
Physlib commit \texttt{3bfc7af} with Lean~\texttt{v4.35.0}.

\section{Mathematical background}\label{sec:background}

Consider a particle of mass $m>0$ at position $\vq\in\mathbb{R}^3$ with momentum $\vp$ in the
potential $V(r)=-k/r$, with Hamiltonian
\[
  H(\vp,\vq)=\frac{\ip{\vp}{\vp}}{2m}-\frac{k}{\|\vq\|}.
\]
The angular momentum is $\vL=\vq\times\vp$ and the Laplace--Runge--Lenz vector is
\[
  \vA=\vp\times\vL-\frac{mk}{\|\vq\|}\,\vq
     =\ip{\vp}{\vp}\,\vq-\ip{\vq}{\vp}\,\vp-\frac{mk}{\|\vq\|}\,\vq ,
\]
where the second form uses $\mathbf a\times(\mathbf b\times\mathbf c)
=\ip{\mathbf a}{\mathbf c}\mathbf b-\ip{\mathbf a}{\mathbf b}\mathbf c$
\cite[\S3.9]{goldstein}. Taking the inner product with $\vq$ and using Lagrange's identity
$\|\vL\|^2=\ip{\vq}{\vq}\ip{\vp}{\vp}-\ip{\vq}{\vp}^2$ gives
\begin{equation}\label{eq:orbit}
  \ip{\vA}{\vq}=\|\vL\|^2-mk\|\vq\| .
\end{equation}
If $\vA$ and $\vL$ are constants of motion and $\theta$ denotes the angle between $\vA$ and
$\vq$, \eqref{eq:orbit} reads $\|\vq\|(mk+\|\vA\|\cos\theta)=\|\vL\|^2$, that is
$\|\vq\|=\frac{\|\vL\|^2/mk}{1+e\cos\theta}$ with eccentricity $e=\|\vA\|/mk$: a conic with
focus at the origin. A direct expansion also gives
\begin{equation}\label{eq:norm}
  \|\vA\|^2=(mk)^2+2mH\|\vL\|^2 ,
\end{equation}
so for $k>0$ and $\vL\neq0$, $e<1$ if and only if $H<0$. Elementary derivations of these facts
can be found in~\cite{goldstein,moriconi2021}.

\section{Formalization}

\subsection{Setting and conventions}\label{sec:defs}

Physlib models phase space for a particle in three dimensions as pairs
$(\vp,\vq)$ of elements of \texttt{EuclideanSpace ℝ (Fin 3)}, and observables as curried functions
\texttt{p → q → ℝ}. The Poisson bracket is
\[
  \{f,g\}=\ip{\nabla_{\vq}f}{\nabla_{\vp}g}-\ip{\nabla_{\vp}f}{\nabla_{\vq}g},
\]
and a central-force system is a structure carrying a mass $m>0$ and a potential
\texttt{V : ℝ → ℝ} smooth on $(0,\infty)$. We follow these conventions exactly so that the
development can be contributed upstream. The Kepler problem and the LRL vector are:

\begin{lean}
noncomputable def kepler (m k : ℝ) (hm : 0 < m) : CentralForce where
  m := m
  V := fun r => -k / r
  m_pos := hm
  contDiffOn_V := contDiffOn_const.div contDiffOn_id
    fun _ hr => (Set.mem_Ioi.mp hr).ne'

noncomputable def lrlVector (m k : ℝ) (p q : EuclideanSpace ℝ (Fin 3)) :
    EuclideanSpace ℝ (Fin 3) :=
  WithLp.toLp 2 (p.ofLp ⨯₃ (angularMomentum p q).ofLp) - (m * k / ‖q‖) • q
\end{lean}

The definition of \texttt{lrlVector} mirrors the textbook formula $\vp\times\vL-(mk/\|\vq\|)\vq$,
reusing Physlib's \texttt{angularMomentum} and Mathlib's cross product on \texttt{Fin 3 → ℝ}. The
\texttt{toLp}/\texttt{ofLp} coercions are needed because the cross product lives on the plain
function type, whereas norms and inner products live on \texttt{EuclideanSpace}.

\begin{remark}
Lean's convention $x/0=0$ makes \texttt{lrlVector} total. A consequence we exploited is that the
orbit equation \eqref{eq:orbit} holds \emph{without} the hypothesis $\vq\neq0$: at $\vq=0$ both
sides vanish. In contrast, \eqref{eq:norm} genuinely requires $\vq\ne0$.
\end{remark}

\subsection{Algebraic identities}\label{sec:algebra}

\begin{lean}
lemma lrlVector_eq (m k : ℝ) (p q : EuclideanSpace ℝ (Fin 3)) :
    lrlVector m k p q = ⟪p, p⟫_ℝ • q - ⟪q, p⟫_ℝ • p - (m * k / ‖q‖) • q

lemma inner_lrlVector_angularMomentum (m k : ℝ) (p q : EuclideanSpace ℝ (Fin 3)) :
    ⟪lrlVector m k p q, angularMomentum p q⟫_ℝ = 0

lemma inner_lrlVector_position (m k : ℝ) (p q : EuclideanSpace ℝ (Fin 3)) :
    ⟪lrlVector m k p q, q⟫_ℝ = angularMomentumSq p q - m * k * ‖q‖

lemma norm_lrlVector_sq (m k : ℝ) (hm : 0 < m) (p : EuclideanSpace ℝ (Fin 3))
    {q : EuclideanSpace ℝ (Fin 3)} (hq : q ≠ 0) :
    ‖lrlVector m k p q‖ ^ 2 =
      (m * k) ^ 2 + 2 * m * (kepler m k hm).hamiltonian p q * angularMomentumSq p q

lemma norm_lrlVector_lt_iff (m k : ℝ) (hm : 0 < m) (hk : 0 < k)
    (p : EuclideanSpace ℝ (Fin 3)) {q : EuclideanSpace ℝ (Fin 3)} (hq : q ≠ 0)
    (hL : angularMomentumSq p q ≠ 0) :
    ‖lrlVector m k p q‖ < m * k ↔ (kepler m k hm).hamiltonian p q < 0
\end{lean}

The triple-product expansion and the orthogonality to $\vL$ are proved componentwise: after
\texttt{ext i; fin\_cases i} the goals become polynomial identities in the six coordinates (with
$mk/\|\vq\|$ as an opaque atom), closed by \texttt{simp} and \texttt{ring}. All later results are
stated in terms of inner products and norms and use \texttt{lrlVector\_eq} to avoid coordinates.
For \eqref{eq:norm} we expand $\ip{\vA}{\vA}$ bilinearly, rewrite with Physlib's Lagrange
identity \texttt{angularMomentumSq\_eq}, clear the denominator $\|\vq\|$ with
\texttt{field\_simp}, and close with \texttt{ring}. One subtlety: rewriting
$\ip{x}{x}=\|x\|^2$ must happen \emph{after} bilinear expansion, otherwise it fires on
$\ip{c\,\vq}{c\,\vq}$ and blocks the expansion. The eccentricity criterion follows from
\eqref{eq:norm} by comparing squares of non-negative reals and a single \texttt{nlinarith} call
with the hint $m\|\vL\|^2>0$.

\subsection{Conservation of the LRL vector}\label{sec:conservation}

For each component $A_i$ we compute the partial gradients explicitly:
\begin{align*}
  \nabla_{\vp}A_i &= 2q_i\,\vp-p_i\,\vq-\ip{\vq}{\vp}\,\mathbf e_i,\\
  \nabla_{\vq}A_i &= \ip{\vp}{\vp}\,\mathbf e_i-p_i\,\vp-\frac{mk}{\|\vq\|}\,\mathbf e_i
     +\frac{mk\,q_i}{\|\vq\|^3}\,\vq \qquad(\vq\neq0).
\end{align*}
In Lean these are \texttt{HasGradientAt} statements, assembled from Physlib's gradient calculus
(\texttt{hasGradientAt\_inner\_self}, \texttt{hasGradientAt\_coord}, \texttt{hasGradientAt\_norm}
and a chain rule \texttt{HasDerivAt.comp\_hasGradientAt}), followed by \texttt{module} or
\texttt{field\_simp} to normalize the resulting linear combination. Using
$V'(r)=k/r^2$ (lemma \texttt{kepler\_deriv\_V}) and Physlib's gradients of the Hamiltonian,
$\nabla_{\vp}H=\vp/m$ and $\nabla_{\vq}H=(k/\|\vq\|^3)\,\vq$, the Poisson bracket evaluates to
\[
  \{A_i,H\}=-k\Big(\frac{p_i}{\|\vq\|}-\frac{q_i\ip{\vq}{\vp}}{\|\vq\|^3}\Big)
  -\frac{k}{\|\vq\|^3}\big(q_i\ip{\vq}{\vp}-p_i\|\vq\|^2\big)=0 .
\]

\begin{lean}
lemma poissonBracket_lrlVector_hamiltonian (m k : ℝ) (hm : 0 < m) (i : Fin 3)
    (p : EuclideanSpace ℝ (Fin 3)) {q : EuclideanSpace ℝ (Fin 3)} (hq : q ≠ 0) :
    poissonBracket (fun p' q' => lrlVector m k p' q' i)
      (kepler m k hm).hamiltonian p q = 0
\end{lean}

The proof is eleven lines: rewrite the bracket via \texttt{poissonBracket\_eq\_of\_hasGradientAt}
with the four gradients, expand inner products, and close with \texttt{field\_simp; ring}.

\subsection{From phase space to trajectories}\label{sec:trajectory}

A vanishing Poisson bracket is the infinitesimal statement of conservation. To obtain Kepler's
first law we need the integrated statement along an actual motion. Physlib provides Hamilton's
equations as \texttt{hamiltonEqOp H p q = 0} for curves \texttt{p q : Time → EuclideanSpace ℝ (Fin 3)},
and the theorem \texttt{deriv\_comp\_eq\_poissonBracket} stating
$\frac{d}{dt}f(\vp(t),\vq(t))=\{f,H\}(\vp(t),\vq(t))$ along solutions. We proved a general lemma:

\begin{lean}
lemma eq_of_poissonBracket_eq_zero {H f : E → E → ℝ} {p q : Time → E}
    (heq : hamiltonEqOp (fun _ => H) p q = 0)
    (hp : Differentiable ℝ p) (hq : Differentiable ℝ q)
    (hf : ∀ t, DifferentiableAt ℝ (Function.uncurry f) (p t, q t))
    (hpb : ∀ t, poissonBracket f H (p t) (q t) = 0) (t s : Time) :
    f (p t) (q t) = f (p s) (q s)
\end{lean}
(here \texttt{E} abbreviates \texttt{EuclideanSpace ℝ (Fin 3)}). Physlib's \texttt{Time} is a
structure wrapping a real number rather than \texttt{ℝ} itself; the proof transports the
derivative to $\mathbb{R}$ with the bridge lemma \texttt{hasDerivAt\_mk\_of\_differentiableAt}
and applies Mathlib's \texttt{is\_const\_of\_deriv\_eq\_zero}. Together with differentiability of
$A_i$ away from the origin and the existing conservation of $\|\vL\|^2$, this yields:

\begin{theorem}[Kepler's first law]
Let $(\vp(t),\vq(t))$ be a differentiable solution of Hamilton's equations for the Kepler
Hamiltonian with $\vq(t)\neq0$ for all $t$. Then for all times $t_0,t$,
\[
  \ip{\vA(t_0)}{\vq(t)}=\|\vL(t_0)\|^2-mk\,\|\vq(t)\|.
\]
\end{theorem}

\begin{lean}
lemma inner_lrlVector_position_of_hamiltonEqOp (m k : ℝ) (hm : 0 < m)
    {p q : Time → EuclideanSpace ℝ (Fin 3)}
    (heq : hamiltonEqOp (fun _ => (kepler m k hm).hamiltonian) p q = 0)
    (hp : Differentiable ℝ p) (hq : Differentiable ℝ q) (hq0 : ∀ t, q t ≠ 0)
    (t₀ t : Time) :
    ⟪lrlVector m k (p t₀) (q t₀), q t⟫_ℝ =
      angularMomentumSq (p t₀) (q t₀) - m * k * ‖q t‖
\end{lean}

Combined with \texttt{norm\_lrlVector\_lt\_iff}, this shows that a bound Kepler motion with
non-zero angular momentum traces an ellipse. Note that we assume, rather than construct, a
solution of Hamilton's equations; existence and uniqueness of solutions are not needed for these
statements and are not addressed.

\section{Discussion}

\paragraph{Size and checking.} The development consists of 2 definitions and 15 lemmas in about
300 lines, including documentation. No proof exceeds a dozen lines. \texttt{\#print axioms} on the
main results reports only \texttt{propext}, \texttt{Classical.choice} and \texttt{Quot.sound}. The
file passes Physlib's module-documentation, alpha-library and style linters.

\paragraph{What was easy.} Once the right normal form (\texttt{lrlVector\_eq}) was available, every
identity of the Kepler problem reduced to \texttt{field\_simp} and \texttt{ring}. Physlib's
existing gradient API and its Poisson-bracket lemma \texttt{poissonBracket\_eq\_of\_hasGradientAt}
made the conservation proof essentially a transcription of the pen-and-paper computation.

\paragraph{What required care.} The friction points were elaboration issues rather than
mathematics: the higher-order unification of \texttt{DifferentiableAt.comp} and of the chain rule
for gradients needed explicit named arguments; \texttt{fun\_prop} could not discharge
differentiability of $\vq\mapsto q_i$ on \texttt{EuclideanSpace} without hand-supplied facts; and
moving between \texttt{Time} and \texttt{ℝ} required Physlib's bridge lemmas.

\paragraph{Future work.} Natural continuations are Kepler's second law (constant areal velocity,
essentially the conservation of $\vL$ along trajectories), Kepler's third law relating the period
to the semi-major axis, the hidden $SO(4)$ symmetry via the Poisson algebra of $\vL$ and
$\vA/\sqrt{-2mH}$ (the classical counterpart of Physlib's quantum commutators), and Bertrand's
theorem. For an upstream contribution to Physlib, the general lemma
\texttt{eq\_of\_poissonBracket\_eq\_zero} should move to the Poisson-bracket module, and the work
should be split into several smaller pull requests as required by the project's guidelines.

\paragraph{Author contributions.} N.L.\ chose the problem, directed the work and reviewed the
results. Claude, an AI system developed by Anthropic, surveyed the existing literature and
libraries, designed the statements, wrote the Lean proofs and drafted the manuscript. All
statements were checked by the Lean kernel. As an AI system cannot take responsibility for a
publication, N.L.\ is accountable for its content.

\bibliographystyle{plain}
\bibliography{refs}

\end{document}
